\subsection{Описание моделируемой системы}

В качестве системы для верификации выбран обмен запросами и ответами
по протоколу Modbus. В рамках работы рассматривается небольшая сеть,
состоящая из одного ведущего устройства Master и трех ведомых устройств
Slave с адресами 1, 2 и 3. Master формирует запрос, все Slave получают
его копию, после чего каждый из них определяет, требуется ли обрабатывать
этот запрос.

Основной операцией является запись одного регистра с кодом функции
\code{0x06}. У каждого Slave предусмотрены два регистра с адресами 0 и 1.
Если запрос содержит адрес конкретного устройства, запись выполняет только
это устройство. После успешной записи оно возвращает ответ, повторяющий
код функции, адрес регистра и записанное значение. Такой порядок обработки
соответствует назначению функции Write Single Register~\cite{modbus}.

Кроме адресной записи, модель предусматривает широковещательный запрос
с адресом 0. Его обрабатывают все три Slave, однако ответ Master
не отправляется. Для проверки отсутствующего получателя используется
адрес 4: такого устройства в сети нет, поэтому запрос остается без ответа.

Перед записью Slave проверяет адрес получателя, CRC, код функции и адрес
регистра. Кадр с неверной CRC игнорируется. При верной CRC неподдерживаемая
функция вызывает исключение с кодом 1. Если функция поддерживается, но
указан несуществующий регистр, возвращается код 2. Исключения отправляются только в ответ на адресные
запросы. На широковещательный запрос Slave не отвечает, даже если в нем
есть ошибка.

\begin{table}[H]
\centering
\caption{Основные параметры модели}
\begin{tabularx}{\linewidth}{@{}p{0.33\linewidth}X@{}}
\toprule
Параметр & Принятое значение \\
\midrule
Состав участников & Один Master и три Slave \\
\addlinespace[3pt]
Адрес получателя & 0 - всем устройствам\newline
1-3 - конкретному Slave\newline
4 - отсутствующему устройству \\
\addlinespace[3pt]
Адрес регистра & 0 или 1\newline
2 - ошибочный адрес \\
\addlinespace[3pt]
Значение для записи & 1 или 2\newline
Начальные значения регистров равны 0 \\
\addlinespace[3pt]
Код функции & 6 - запись одного регистра\newline
3 - неподдерживаемая функция \\
\addlinespace[3pt]
Корректность CRC & Логический признак \code{crc_ok} \\
\addlinespace[3pt]
Число активных запросов & Одна транзакция за раз \\
\bottomrule
\end{tabularx}
\end{table}

Master повторяет обмен в бесконечном цикле. Параметры очередного запроса
выбираются недетерминированно, поэтому верификатор рассматривает различные
запросы и допустимые порядки выполнения процессов. Адрес получателя,
регистр, значение, функция и корректность CRC выбираются независимо.
Поэтому запрос может быть корректным или содержать сразу несколько ошибок.
Например, Master может отправить функцию \code{FC03} с адресом регистра 2
и неверной CRC. Всего предусмотрено 120 сочетаний параметров запроса:
пять адресов получателя, три адреса регистра, два значения, две функции
и два состояния CRC.

\subsection{Взаимодействие процессов и принятые предположения}

На рисунке~\ref{fig:exchange} приведен пример адресного обмена со Slave 2.
Общий запрос доставляется каждому устройству через отдельный канал.
Slave 1 и Slave 3 проверяют адрес и завершают обработку без изменения
регистров. Slave 2 выполняет запись и помещает ответ в общий канал ответов.
Уведомление о завершении отправляет каждый из трех обработчиков.

\begin{figure}[!ht]
\centering
\includegraphics[width=\linewidth,height=0.60\textheight,keepaspectratio]{exchange.png}
\caption{Один из порядков взаимодействия процессов при запросе к Slave 2}
\label{fig:exchange}
\end{figure}

Каналы запросов имеют емкость в одно сообщение. Это позволяет каждому
Slave независимо получить свою копию одного и того же кадра. Общий канал
ответов также рассчитан на одно сообщение, поскольку адресным получателем
может быть только одно устройство. Служебный канал завершения вмещает три
уведомления. Эти уведомления нужны для согласования шагов модели и
не являются сообщениями протокола Modbus.

Результат транзакции фиксируется после получения всех трех уведомлений.
При адресном запросе наличие сообщения в канале ответов означает получение
ответа. Если канал пуст, Master фиксирует тайм-аут. При широковещательном
запросе Master сразу фиксирует завершение без ответа. Таким образом,
тайм-аут описывает логический результат обмена; числовой интервал ожидания
в модели не задается.

\needspace{6\baselineskip}
Передача кадра представлена доставкой полей сообщения. Корректность CRC
задается логическим признаком. Это позволяет сосредоточиться на выборе получателя,
условиях записи, содержимом ответа и завершении транзакции. Значения
регистров хранятся в переменных типа \code{byte}, а генератор выбирает
только 1 и 2. Проверка показывает, как система обрабатывает запросы
с выбранными значениями параметров.

Верифицируемые свойства должны показать, что корректные запросы выполняют
нужную запись, ошибочные запросы сохраняют память, посторонние устройства
и регистры не изменяются, ответы формируются правильно и начатый обмен
обязательно достигает результата. Для этого далее рассматривается
реализация модели и средства наблюдения за каждой транзакцией.
